% template.tex: a starter document for the aurea style.
%
% Put aurea.sty next to this file (or in TEXMFHOME/tex/latex/aurea/) and
% build with:  latexmk -pdf template.tex   (lualatex and xelatex also work)
%
% The class sets the size and the paper (aurea's golden block is designed for
% A4 portrait); aurea does the rest. For a DRAFT watermark, uncomment `draft'
% below. Putting `draft' in the class options also gives the watermark, but
% then graphicx prints figures as empty boxes and hyperref drops links and
% the PDF metadata.
\documentclass[12pt,a4paper]{article}

% Every package option. The first five show their default value; delete what
% you do not change. The last three are off unless you uncomment them.
\usepackage[
  language  = italian,  % babel languages, main one last: {latin,italian}
  figures   = osf,      % osf (old-style) or lining
  nonumbers = false,    % true: unnumbered chapters and sections
  geometry  = {},       % extra geometry keys, e.g. {top=40mm}
  urlstyle  = tt,       % \url font: tt, rm, sf or same
  % draft,              % DRAFT watermark (figures and links stay normal)
  % watermark = BOZZA,  % watermark with this text (DRAFT when left empty)
  % final,              % no watermark, even with the `draft' class option
]{aurea}

% Title page and PDF metadata in one place. Values that contain commas go in
% braces, e.g. authornote = {Parrocchia S. Lorenzo da Brindisi, Francavilla
% Fontana}; without them the part after the comma is read as an unknown key.
% \title, \author, \date, \subtitle and \authornote also work alone.
\aureasetup{
  title      = Titolo del documento,
  subtitle   = Sottotitolo in corpo minore,
  author     = Nome dell'autore,
  authornote = {qualifica, luogo o occasione},
  date       = 24 settembre 2026,
  subject    = {Una frase che descrive il documento.},
  keywords   = {parola chiave, altra parola chiave},
}

% Per-document macros go here (in ce-Speranza-e-speranza, for example,
% \puntata stays here and does not move into the package).

\begin{document}

% The title page gets no page number: aurea adds \thispagestyle{empty}.
\maketitle

% An epigraph, if there is one, stays on the title page.
\begin{citazione}[Rm 4,18]
Egli ebbe fede sperando contro ogni speranza e così divenne padre di molti
popoli.
\end{citazione}

% The abstract on a page of its own: sommario is \clearpage, abstract (babel
% titles it "Sommario"), \clearpage. The page keeps its number, as in the
% four projects.
\begin{sommario}
Un paragrafo che riassume il documento, su una pagina a sé.
\end{sommario}

% A lead artwork, if there is one. Budget room for the caption in the height,
% so that the float does not grow taller than the page.
\begin{figure}[p]
  \centering
  \includegraphics[width=\linewidth,height=0.8\textheight,keepaspectratio]
    {example-image}
  \caption[Titolo breve]{\emph{Titolo dell'opera}, didascalia estesa.}
\end{figure}

\clearpage

\section{Primo capitolo}

Il testo corre con il rientro di prima riga e senza spazio fra i paragrafi.
Le parole dette vanno fra caporali in corsivo, \detto{come questa}, anche
annidate: \detto{disse \detto{seguimi} e se ne andò}. Un riferimento biblico
letterale va in nota con \verb|\rif|\rif{Mt 9,9}, un'allusione con
\verb|\cfr|\cfr{Gv 18,17.25}. Le fonti in rete vanno in nota con
\verb|\url|.\footnote{\url{https://www.example.org/}}

\begin{citazione}
Una citazione in blocco senza fonte: il testo è in corsivo, senza
\verb|\emph|.
\end{citazione}

\scenebreak

Dopo la riga aurea riprende il racconto. Una lettera si chiude con la firma,
un paragrafo a sé in corsivo:

\begin{quote}
\emph{«Ti benedico».}

\firma{Padre Pio, cappuccino}
\end{quote}

\section{Secondo capitolo}

Il nome usato nei primi progetti, \verb|\sectionbreak|, resta come
sinonimo:

\sectionbreak

Fine.

% The table of contents ("Indice") at the back, on a page of its own.
\indice

\end{document}
